% ===== 原文 pp.45–46（§5.3 Langlands 对偶）=====
% [原文 p.45]

\subsection{Langlands 对偶}

模空间 $\cM_{G_{\C}}$ 的各类子空间之间的最后一种对应关系来自 Langlands 对偶。考虑膜（brane）的主要理由之一，在于它与镜像对称以及几何 Langlands 纲领之间的联系。设 ${}^{L}G_{\C}$ 为 $G_{\C}$ 的 Langlands 对偶群，则 $G_{\C}$ 与 ${}^{L}G_{\C}$ 的不变多项式之间存在对应，从而给出 Hitchin 底的一个等同 $\cA_{G_{\C}} \simeq \cA_{{}^{L}G_{\C}}$。于是两个模空间 $\cM_{G_{\C}}$ 与 $\cM_{{}^{L}G_{\C}}$ 都是同一底空间上的环面纤维化，且其非奇异纤维互为对偶阿贝尔簇 [40]。Kapustin 与 Witten 借助 S-对偶给出了这一点的物理解释，并以此作为他们研究几何 Langlands 纲领的基础 [69]。

% [Tutor Note: p.45 原文 "an identification ... of the Hitchin basis" 中的 "basis" 当为 "base" 之误：$\cA_{G_{\C}}$ 是由不变多项式构成的 Hitchin 底（全书他处均作 Hitchin base），并非某个「基」。正文按原文语气译作「Hitchin 底」。又，同句 "the moduli space of $G_{\C}$-Higgs bundle" 中 "bundle" 为单数，疑漏字母 s；中文无单复数之别，此处仅记。]

在这一做法中，起关键作用的是各种类型的膜以及它们在镜像对称下的变换。Langlands 对偶按如下方式交换膜的类型：
\[
\begin{gathered}
(B,B,B) \leftrightarrow (B,A,A),\\
(A,B,A) \leftrightarrow (A,B,A),\\
(A,A,B) \leftrightarrow (A,A,B),
\end{gathered}
\]
但确切的对应关系尚不清楚。此外，借助 Hitchin 纤维化研究 Higgs 丛的膜，可以考虑 Fourier--Mukai 型变换，从而用对偶纤维的子空间来得到对偶；关于第 4 章构造的膜与 Langlands 对偶，有几点可以说：

\begin{itemize}[leftmargin=1.4em,itemsep=1pt,topsep=1pt]
\item \textbf{3-流形膜。}回忆 [17] 中有一个自然的 $(A,B,A)$-膜，它是 $i_2(\bar\partial_A,\Phi) = (f^*(\partial_A), f^*(\Phi^*))$ 的不动点，其中 $f:\Sigma\to\Sigma$ 是 $\Sigma$ 上的实结构，$\rho$ 是 $G_{\C}$ 的紧实形式。借助 Langlands 对偶群的紧形式 ${}^{L}\rho$，这些空间在 $\cM_{{}^{L}G_{\C}}$ 上有相应的对合，由下式给出：
\[
{}^{L}i_2(\bar\partial_A,\Phi) = (f^*(\partial_A), f^*(\Phi^*)) = (f^*({}^{L}\rho(\bar\partial_A)), -f^*({}^{L}\rho(\Phi))),
\]
其不动点集在 [17] 中被猜想为对偶膜（细节亦见 [90, 猜想 1]）。
% [Tutor Note: p.45 此式中间一项 $(f^*(\partial_A),f^*(\Phi^*))$ 与上文 $i_2$ 在 $\cM_{G_{\C}}$ 上的定义逐字相同，而左端 ${}^{L}i_2$ 是对偶空间 $\cM_{{}^{L}G_{\C}}$ 上的对合，两者不应相等；按 [90, Conjecture 1] 的语境，对偶对合应直接由最右端给出。疑中间一项系原文误植（或应在其中插入 ${}^{L}\rho$）。正文照原样排，建议对照 [17]、[90] 原式复核。]
\item \textbf{伪实 Higgs 丛。}如 [17] 所示，$\cM_{G_{\C}}$ 中的伪实 Higgs 丛构成一个 $(A,A,B)$-膜，并可看作对合 $i_3(\bar\partial_A,\Phi) = (f^*\sigma(\bar\partial_A), f^*\sigma(\Phi))$ 的不动点集，其中 $\sigma$ 是复李群 $G_{\C}$ 的一个实形式。此时 [17] 猜想对偶膜是如下定义的对应对合的不动点集：
\[
{}^{L}i_3(\bar\partial_A,\Phi) = (f^*\hat\sigma(\bar\partial_A), f^*\hat\sigma(\Phi)),
\]
其中 $\hat\sigma$ 是 ${}^{L}G_{\C}$ 上的实结构，并猜想它是 $\sigma$ 的 Nadler 对偶（细节见下文）。
% [原文 p.46]
\item \textbf{超多边形空间。}以箭图方式出现的膜可能带有群结构，因此可以对它们考虑 Langlands 对偶。关于由箭图表示产生的膜之间的对偶，一些一般性结果见 [66]；[17] 意义下的超多边形膜的几何性质则见 [85]。
\end{itemize}

在 Higgs 丛的膜的 Langlands 对偶中，最有趣的情形之一是考虑实 Higgs 丛构成的 $(B,A,A)$-膜 $\cM_G\subset\cM_{G_{\C}}$，它对应于紧黎曼面 $\Sigma$ 的曲面群的实表示。回忆这些膜可以通过 $G_{\C}$ 的实形式 $G$ 的 Cartan 对合 $\theta$ 来看，即作为对合
\[
i_1(\bar\partial_A,\Phi) = (\theta(\bar\partial_A), -\theta(\Phi))
\]
的不动点集。

对偶的 $\cM^{L}_{G}\subset{}^{L}\cM_{{}^{L}G_{\C}}$ 必须是 $(B,B,B)$ 型，即一个关于 $I,J,K$ 全纯的子流形，其上承载超全纯层（hyperholomorphic sheaf）。据 [17]，所提出的对偶膜应取如下形式（细节亦见 [90, 猜想 2]）：

\begin{conj}{5.13}
与 $i_1$ 的不动点集对偶的膜，其支撑是模空间 $\cM_{\check H}\subset\cM_{{}^{L}G_{\C}}$，即 $\check H$-Higgs 丛的模空间，其中 $\check H$ 是对应于 [76, 表 1] 中出现的李代数 $\check\fh$ 的群。
\end{conj}
% [Tutor Note: p.46 猜想 5.13 中的 $\check H$ 依所取实形式 $G$ 而定：$i_1$ 由 $G$ 的 Cartan 对合 $\theta$ 定义，而对偶群的候选 $\check H$（见 [76, Table 1]）按 $G$ 的实形式分情形列出。原文在猜想中未显式写出对 $G$ 的依赖，应理解为：对每个实形式 $G$，对偶膜的支撑由表中相应的 $\check H$ 给出。]

尽管上述猜想尚未得到证明，但各个方向上已有大量工作取得进展。下面我们提及其中一些工作：

\begin{itemize}[leftmargin=1.4em,itemsep=1pt,topsep=1pt]
\item \textbf{带非阿贝尔数据的膜。}如前所述，$\Sl(m,\HH)$、$\SO(2n,\HH)$ 与 $\Sp(2m,2m)$-Higgs 丛的 $(B,A,A)$-膜与 Hitchin 纤维化的交被证明是非阿贝尔的 [64]。这些膜的对偶不久前在 Branco 的博士论文（DPhil thesis）中作了研究：作者先描述了包含这些膜的整个奇异纤维的几何，然后才得到对偶空间 [30]。
% [Tutor Note: p.46 所谓与 Hitchin 纤维化的「非阿贝尔交」，指这些膜与 Hitchin 纤维的交不是有限个挠点（0 维），而是具有正维数、由非阿贝尔对象构成的子簇；这与下文「分裂实形式」情形中交为 2-挠点集（0 维）恰成对照。[64] 正是以这一现象刻画 $\Sl(m,\HH)$ 等膜。]
\item \textbf{拟分裂实形式。}$\U(m,m)$-Higgs 丛的 $(B,A,A)$-膜与 Hitchin 纤维化正则纤维的交已在 [91] 中研究，通过这一描述，[65] 给出了一个与猜想 5.13 相符的对偶膜的候选。此外，对 $m=1$，最近 [99] 利用上同调类与拓扑不变量为该猜想提供了进一步支持。
\item \textbf{分裂实形式。}对 $G_{\C}$ 的分裂实形式 $G$，[89, 定理 4.12] 证明由对合 $i_1$ 得到的 $(B,A,A)$-膜在 Hitchin 底上方以 2-挠点的形式出现，因而在每根纤维中是 0 维空间。对膜的显式描述（以及大多数经典分裂实群的单值性作用 [12,16,88,89]）应当为研究镜像对称提供一个良好的起点，值得注意的是它与 Hitchin 在 [65] 中所研究情形的相似性。
% [Tutor Note: p.46 「在 Hitchin 底上方以 2-挠点的形式出现」意为：该膜在 Hitchin 映射下映满整个底，而它与每根正则纤维（一个阿贝尔簇）的交是纤维中的 2-挠点；阿贝尔簇的 2-挠点群为 $(\Z/2\Z)^{2d}$（$d$ 为纤维维数），故交为有限集、0 维。可据此核对本句末尾 "hence a 0-dimensional space in each fibre"。]
\end{itemize}
